
%%%%%%%%%%%%%%%%%%%%%%%%%%%%%%%%%%%%%%%%%%%%%%%%%%
%% 第十章 ξ 函数——整函数化与 Hadamard 乘积
%%%%%%%%%%%%%%%%%%%%%%%%%%%%%%%%%%%%%%%%%%%%%%%%%%

\chapter{\texorpdfstring{$\xi$}{xi(s)} 函数——整函数化与 Hadamard 乘积}\label{ch:xi_function}

\section*{引言}\label{sec:intro_xi_function}

$\zeta(s)$ 的函数方程体现关于 $\operatorname{Re}(s)=1/2$ 的对称，但 $\zeta$ 在 $s=1$ 有极点，且函数方程中含 $\Gamma(1-s)$、$\sin(\pi s/2)$ 等因子。定义
\[
\xi(s)=\tfrac12 s(s-1)\pi^{-s/2}\Gamma(s/2)\zeta(s)
\]
后，极点被消去，$\xi$ 为整函数，且 $\xi(s)=\xi(1-s)$。

本章分三步展开：
\begin{enumerate}
  \item \textbf{整函数化}：定义 $\xi$，论证其为整函数并导出 $\xi(s)=\xi(1-s)$，说明零点与非平凡零点的对应，并给出曲面示意；
  \item \textbf{Hadamard 乘积}：说明阶与基本因子 $E(z,1)$，写出乘积公式与常数 $b_0$，讨论黎曼猜想（RH）下的简化形式，并导出黎曼原稿中的 $t$-乘积；
  \item \textbf{对数导数}：由定义与乘积给出 $\xi'/\xi$，换算为 $\zeta'/\zeta$，供第~\ref{ch:riemann_explicit_formula}~章使用。
\end{enumerate}
几何分布见第~\ref{ch:zero_distribution}~章；与原稿的对照见第~\ref{ch:riemann_original_interpretation}~章。

\section{\texorpdfstring{$\xi$}{xi(s)} 函数的定义与基本性质}\label{sec:xi-definition}
%%%%%%%%%%%%%%%%%%%%%%%%%%%%%%%%%%%%%%%%%%%%%%%%%%

\subsection{构造的动机}

第~\ref{ch:riemann_zeta_continuation}~章已建立 $\zeta$ 的全平面亚纯延拓与非对称函数方程
\[
\zeta(s) = 2^s \pi^{s-1} \sin\!\left(\frac{\pi s}{2}\right) \Gamma(1-s)\, \zeta(1-s).
\]
该方程体现关于 $\operatorname{Re}(s)=\frac12$ 的对称，但形式上有两点不便：
\begin{itemize}
    \item $\zeta$ 在 $s=1$ 有单极点，不是整函数；
    \item 右端含 $\Gamma(1-s)$、$\sin(\pi s/2)$ 等因子，对称 $s\mapsto 1-s$ 不够直接。
\end{itemize}
定义辅助函数 $\xi(s)$，用因子消去极点与上述附加结构，使对称写成 $\xi(s)=\xi(1-s)$。

\subsection{\texorpdfstring{$\xi(s)$}{xi(s)} 的精确定义}

\begin{definition}{\texorpdfstring{Riemann $\xi$ 函数}{Riemann xi 函数}}{}
定义 Riemann $\xi$ 函数为：
\begin{equation}
    \xi(s) \;=\; \frac{1}{2}\, s(s-1)\, \pi^{-s/2}\, \Gamma\!\left(\frac{s}{2}\right)\, \zeta(s).
    \label{eq:xi-definition}
\end{equation}
\end{definition}

该定义中各因子的作用如下：
\begin{itemize}
    \item \textbf{因子 $\pi^{-s/2}\,\Gamma(s/2)$}：$\zeta$ 函数方程中的完备化因子；与 $\zeta(s)$ 合并得
    $\xi^*(s)=\pi^{-s/2}\Gamma(s/2)\zeta(s)$，满足 $\xi^*(s)=\xi^*(1-s)$。
    \item \textbf{因子 $\frac{1}{2}s(s-1)$}：$s=1$ 处 $(s-1)$ 消去 $\zeta$ 的单极点；
    $s=0$ 处 $s$ 消去 $\Gamma(s/2)$ 的单极点。
    在 $s=-2,-4,\ldots$ 处，$\Gamma(s/2)$ 的极点与 $\zeta$ 的平凡零点相消，故 $\xi$ 无极点。
    前置系数 $\frac12$ 使 $\xi(0)=\xi(1)=\frac12$（正规化）。
\end{itemize}

\subsection{整函数性：极点与平凡零点的相消}
\label{sec:xi-entire-proof}

由定义~\eqref{eq:xi-definition}，只需说明右端在 $\mathbb{C}$ 上可解析延拓为全纯函数（各潜在奇点处的奇性相互抵消）。
以下逐点核对；所用事实分别见第~\ref{ch:gamma_function}~章与第~\ref{ch:riemann_zeta_continuation}~章。

\paragraph{$s=1$}
$\zeta$ 在 $s=1$ 有一阶极点，且 $\operatorname{Res}_{s=1}\zeta(s)=1$。
因子 $(s-1)$ 恰好消去该极点，故
\[
  (s-1)\zeta(s)
\]
在 $s=1$ 处全纯且取值 $1$。其余因子 $\frac12 s\,\pi^{-s/2}\Gamma(s/2)$ 在 $s=1$ 全纯非零，因而 $\xi$ 在 $s=1$ 全纯。

\paragraph{$s=0$}
$\Gamma(s/2)$ 在 $s=0$ 有一阶极点。
由 $\Gamma$ 在 $0$ 附近的 Laurent 展开（或 $\Gamma(z)\sim 1/z$ 当 $z\to 0$），
$\Gamma(s/2)$ 的主部与 $s$ 同阶；因子 $s$ 消去该极点。
$\zeta(0)=-\frac12\neq 0$，$(s-1)$ 与 $\pi^{-s/2}$ 在 $s=0$ 全纯非零，故 $\xi$ 在 $s=0$ 全纯。

\paragraph{平凡零点 $s=-2,-4,\ldots$}
设 $s=-2n$（$n\in\mathbb{N}$）。
此时 $\zeta(-2n)=0$（平凡零点，第~\ref{ch:riemann_zeta_continuation}~章），
而 $\Gamma(s/2)=\Gamma(-n)$ 有极点。
因 $\zeta$ 在 $s=-2n$ 的零点至少为一阶，与 $\Gamma(s/2)$ 的一阶极点相消后，
乘积 $\Gamma(s/2)\zeta(s)$ 在 $s=-2n$ 全纯；
再乘以多项式因子 $\frac12 s(s-1)\pi^{-s/2}$ 仍全纯。
（更高阶零点只会使相消后仍全纯，不引入新极点。）

\paragraph{其余点}
$\pi^{-s/2}$ 为整函数；$s(s-1)$ 为多项式。
$\zeta$ 在 $\mathbb{C}\setminus\{1\}$ 全纯，$\Gamma$ 在 $\mathbb{C}\setminus\{0,-1,-2,\ldots\}$ 全纯，
而上述各点的奇性均已处理。故 $\xi$ 在全平面全纯，即 $\xi$ 为整函数。

\begin{proposition}{\texorpdfstring{$\xi$}{xi} 为整函数}{}
\label{prop:xi-entire}
由定义~\eqref{eq:xi-definition} 与第~\ref{ch:gamma_function}、\ref{ch:riemann_zeta_continuation}~章的延拓结果，
$\xi(s)$ 在 $\mathbb{C}$ 上全纯。
\end{proposition}

\subsection{函数方程（对称性）}
\label{sec:xi-symmetry-sketch}

记完备因子
\[
  \xi^*(s):=\pi^{-s/2}\Gamma\!\left(\frac s2\right)\zeta(s).
\]
第~\ref{ch:riemann_zeta_continuation}~章的函数方程可写成对称形式
\[
  \xi^*(s)=\xi^*(1-s)
\]
（与非对称形式
$\zeta(s)=2^s\pi^{s-1}\sin(\pi s/2)\Gamma(1-s)\zeta(1-s)$
经 $\Gamma$ 的倍元/反射公式等价；推导见该章）。
定义~\eqref{eq:xi-definition} 即
\[
  \xi(s)=\frac12 s(s-1)\,\xi^*(s).
\]
因 $s(s-1)=(1-s)\bigl((1-s)-1\bigr)$，有
\[
  \frac12 s(s-1)\,\xi^*(s)
  =\frac12 (1-s)\bigl((1-s)-1\bigr)\,\xi^*(1-s),
\]
即 $\xi(s)=\xi(1-s)$。

\begin{proposition}{函数方程}{}
\label{prop:xi-fe}
对一切 $s\in\mathbb{C}$，
\begin{equation}
        \xi(s) = \xi(1-s).
        \label{eq:xi-functional-equation}
\end{equation}
\end{proposition}

这是关于临界线 $\Re(s) = \tfrac{1}{2}$ 的镜像对称，形式比 $\zeta$ 的函数方程简洁。
若 $\xi(\rho)=0$，则 $\xi(1-\rho)=\xi(\bar\rho)=\xi(1-\bar\rho)=0$
（后两项还用到 $\xi(\bar s)=\overline{\xi(s)}$）。

\subsection{零点的精确对应}

由定义~\eqref{eq:xi-definition}，$\xi$ 的零点恰好是 $\zeta$ 的非平凡零点（临界带 $0 < \Re(s) < 1$ 内的零点）。
平凡零点被 $\Gamma(s/2)$ 的极点抵消，$s=0,1$ 的极点亦已消去。因此
\[
\xi(\rho) = 0 \;\iff\; \zeta(\rho) = 0 \;\text{ 且 }\; 0 < \Re(\rho) < 1.
\]
\textbf{RH 等价于：$\xi$ 的全部零点满足 $\Re(\rho) = \tfrac{1}{2}$}。

\subsection{\texorpdfstring{$\xi(s)$}{xi(s)} 的几何面貌：实部、虚部与模曲面}

下面给出 $\xi$ 的实部、虚部与模曲面示意。
图为示意，\textbf{不代替}命题~\ref{prop:xi-entire}、\ref{prop:xi-fe} 的证明。
零迹线交点与 RH 的几何表述见第~\ref{ch:zero_distribution}~章。

如图~\ref{fig:xi-real-imag}，实部与虚部以临界线 $\Re(s)=\tfrac12$ 为对称轴（实部偶对称、虚部奇对称，由 $\xi(s)=\xi(1-s)$ 与共轭对称）。
零点处两曲面同时过 $z=0$；局域定位宜用零迹线交点，而非单张曲面。

\begin{figure}[htbp]
\centering
\includegraphics[width=1.1\textwidth]{全书图形/xi实虚曲面联排.pdf}
\caption{$\xi$ 的实部曲面（左）与虚部曲面（右）。
关于 $\Re(s)=\tfrac12$ 分别呈对称与反对称；可视范围内光滑无极点。图为示意。}
\label{fig:xi-real-imag}
\end{figure}

图~\ref{fig:xi-modulus} 为模曲面 $z=|\xi(s)|$：关于临界线与实轴镜面对称，$|\xi|$ 随 $|s|$ 增大而迅速增长（与一阶整函数增长一致）。
局域零点在宏观模曲面上不易辨识。

\begin{figure}[htbp]
\centering
\includegraphics[width=1.0\textwidth]{全书图形/xi模曲面.pdf}
\caption{$\xi$ 的模曲面 $z=|\xi(s)|$。
模非负；关于 $\Re(s)=\tfrac12$ 与实轴 $\Im(s)=0$ 镜面对称。图为示意。}
\label{fig:xi-modulus}
\end{figure}

\section{\texorpdfstring{$\xi(s)$}{xi(s)} 的 Hadamard 无穷乘积与对数导数}\label{sec:xi-hadamard-product}

\subsection{从多项式到 Weierstrass--Hadamard 乘积}

设 $P(z)$ 是 $n$ 次多项式，根为 $\rho_1,\ldots,\rho_n$（计重数），则
\[
P(z) = P(0)\prod_{k=1}^n\left(1-\frac{z}{\rho_k}\right) \tag{1}
\]
Weierstrass 将此推广至整函数，引入基本因子
\[
E(z, p) = (1-z)\exp\!\left(z + \frac{z^2}{2} + \cdots + \frac{z^p}{p}\right), \qquad p \geq 1; \quad E(z,0) = 1-z,
\tag{2}
\]
满足 $|z|\le\frac12$ 时 $|1-E(z,p)|\le C|z|^{p+1}$。
适当选取 $p_n$ 可使
\begin{equation}
f(z) = e^{g(z)} \prod_{n=1}^\infty E\left(\frac{z}{a_n}, p_n\right)
\label{eq:weierstrass-general}
\tag{3}
\end{equation}
在任意紧集上一致收敛。
Hadamard 进一步指出：有限阶整函数可用统一的 $p$，且 $g$ 为次数不超过该阶的多项式。
本章主线是 $\xi$ 的 Hadamard 乘积（按\textbf{零点}分解）。
$\Gamma$ 的 Weierstrass 乘积（按极点 / $1/\Gamma$ 的零点编码）正文主线不引用；
对照证明例见附录~\ref{app:weierstrass}（选读）。
本章采用有限阶整函数的标准结论，并确定 $\xi$ 上应取 $p=1$。

\subsection{增长阶、亏格与基本因子的选取}
\label{sec:xi-order-genus}

整函数 $f$ 的\textbf{阶}为
\[
  \rho_f
  :=\limsup_{r\to\infty}
  \frac{\log\log M(r)}{\log r},
  \qquad
  M(r)=\max_{|z|=r}|f(z)|.
\]
阶有限时可用统一基本因子 $E(z,p)$；亏格由零点收敛指数与阶决定
（一般理论见 Titchmarsh 等；$\Gamma$ 乘积作为 Weierstrass 证明例见附录~\ref{app:weierstrass}）。

\paragraph{$\xi$ 的阶为 $1$}
Stirling 公式（第~\ref{ch:gamma_function}~章）在垂直带内给出
\[
  \log\bigl|\Gamma(\sigma+it)\bigr|
  =\bigl(\sigma-\tfrac12\bigr)\log|t|-|t|\frac{\pi}{2}+O(1).
\]
$\xi$ 中 $\Gamma(s/2)$ 主导指数增长，$\zeta(\sigma+it)$ 在固定带内至多多项式型界
（凸性 / Phragm\'en--Lindel\"of 与函数方程；完整估计见标准教材）。
故 $\log\log M_\xi(r)\sim\log r$，即 $\xi$ 为\textbf{一阶}整函数；
$N(T)\asymp T\log T$（第~\ref{ch:zero_distribution}~章）与此相容。

\paragraph{为何采用 $E(z,1)=(1-z)e^{z}$}
对零点 $\{\rho\}$，$\sum 1/|\rho|$ 发散而 $\sum 1/|\rho|^{1+\varepsilon}$ 收敛，收敛指数为 $1$，故取
\[
  E\!\left(\frac{s}{\rho},1\right)
  =\Bigl(1-\frac{s}{\rho}\Bigr)e^{s/\rho}.
\]
省略 $e^{s/\rho}$ 时乘积一般发散；保留后因子为 $1+O(|s|^2/|\rho|^2)$，
在按 $|\operatorname{Im}\rho|$ 的对称截断下于紧集一致收敛。
指数因子在阶 $1$ 时至多一次，故
\[
  \xi(s)=\xi(0)\,e^{b_0 s}
  \prod_{\rho}
  \Bigl(1-\frac{s}{\rho}\Bigr)e^{s/\rho}.
\]

\begin{proposition}{\texorpdfstring{$\xi$}{xi} 为一阶整函数}{}
\label{prop:xi-order-one}
由 Stirling 公式对 $\Gamma(s/2)$ 的控制与垂直带内 $\zeta$ 的多项式型界，$\xi$ 为一阶整函数；
Hadamard 乘积取 $E(z,1)$，指数因子为一次多项式。
完整增长理论见 Titchmarsh 等。
\end{proposition}

\subsection{Hadamard 乘积与对数导数}
\label{sec:xi-to-zeta-logderiv}

由命题~\ref{prop:xi-order-one}，在按 $|\operatorname{Im}\rho|$ 对称截断的约定下，
\begin{equation}
\boxed{
\xi(s) = \xi(0) e^{b_0 s} \prod_{\rho} \left(1 - \frac{s}{\rho}\right) e^{s/\rho}
}
\label{eq:hadamard-product-ch9}
\tag{4}
\end{equation}
其中 $\rho$ 跑遍非平凡零点（计重数），$\xi(0)=\frac12$，且
\[
  b_0
  \;=\;
  -\frac{\gamma}{2}-1+\frac12\log(4\pi)
  \;=\;
  \log 2-1-\frac{\gamma}{2}+\frac12\log\pi
\]
（$\gamma$ 为 Euler--Mascheroni 常数）。
与 (1) 对照：$\xi(0)$ 对应 $P(0)$，因子 $(1-s/\rho)e^{s/\rho}$ 即 $E(s/\rho,1)$。

\paragraph{两种 $\xi'/\xi$}
对~\eqref{eq:hadamard-product-ch9} 取对数导数（$s=0$ 邻域，避开零点）：
\begin{equation}
\frac{\xi'(s)}{\xi(s)} = b_0 + \sum_{\rho} \left( \frac{1}{s - \rho} + \frac{1}{\rho} \right).
\label{eq:xi-log-derivative}
\tag{5}
\end{equation}
于是 $b_0=\xi'(0)/\xi(0)$。
另一方面，由定义~\eqref{eq:xi-definition} 的乘积法则（$\operatorname{Re}s$ 充分大处微分，再解析延拓）：
\begin{equation}
  \frac{\xi'(s)}{\xi(s)}
  =
  \frac{1}{s}+\frac{1}{s-1}
  -\frac12\log\pi
  +\frac12\frac{\Gamma'(s/2)}{\Gamma(s/2)}
  +\frac{\zeta'(s)}{\zeta(s)}.
  \label{eq:xi-logderiv-from-def}
\end{equation}
令 $s\to 0$，用 $\zeta(0)=-\frac12$、$\zeta'(0)/\zeta(0)=\log(2\pi)$ 及
$\Gamma'(z)/\Gamma(z)=-1/z-\gamma+O(z)$，即得上述 $b_0$ 闭式。

由 $\xi(s)=\xi(1-s)$ 得 $\xi'(s)/\xi(s)=-\xi'(1-s)/\xi(1-s)$，故 $\xi'(\frac12)=0$。
对称截断下还有 $b_0+\sum_{\rho}1/\rho=0$，从而 RH 下
$b_0+\sum_{\tau>0}1/(\tfrac14+\tau^2)=0$，用于吸收下文 (7) 的全局指数。

\paragraph{换算到 $\zeta'/\zeta$}
显式公式使用的是 $-\zeta'/\zeta$。由~\eqref{eq:xi-logderiv-from-def}，
\begin{equation}
  \frac{\zeta'(s)}{\zeta(s)}
  =
  \frac{\xi'(s)}{\xi(s)}
  -\frac{1}{s}-\frac{1}{s-1}
  +\frac12\log\pi
  -\frac12\frac{\Gamma'(s/2)}{\Gamma(s/2)}.
  \label{eq:zeta-from-xi-logderiv}
\end{equation}
再代入~\eqref{eq:xi-log-derivative}，
\begin{equation}
  \frac{\zeta'(s)}{\zeta(s)}
  =
  b_0
  +\sum_{\rho}\Bigl(\frac{1}{s-\rho}+\frac{1}{\rho}\Bigr)
  -\frac{1}{s}-\frac{1}{s-1}
  +\frac12\log\pi
  -\frac12\frac{\Gamma'(s/2)}{\Gamma(s/2)}.
  \label{eq:zeta-logderiv-hadamard}
\end{equation}
在 $\operatorname{Re}s>1$ 时左端等于 $-\sum_{n\ge 1}\Lambda(n)n^{-s}$。
第~\ref{ch:riemann_explicit_formula}~章将~\eqref{eq:zeta-logderiv-hadamard} 代入 Perron 积分公式并取留数。

\subsection{零点配对、RH 下的形式与原稿 $t$-乘积}
\label{sec:hadamard-to-riemann-product}

利用 $\xi(s)=\xi(1-s)$，零点成 $\rho$ 与 $1-\rho$ 配对。一般位置下两项乘积为
\[
\left(1 - \frac{s}{\rho}\right) e^{s/\rho} \cdot \left(1 - \frac{s}{1-\rho}\right) e^{s/(1-\rho)};
\]
仅当 $\operatorname{Re}\rho=\frac12$ 时退化为关于 $s-\frac12$ 的实二次型。

若 RH 成立，$\rho=\frac12\pm i\tau$（$\tau>0$），配对给出
\[
\left(1 - \frac{s}{\frac{1}{2} + i\tau}\right)
\left(1 - \frac{s}{\frac{1}{2} - i\tau}\right)
\cdot
\exp\left(\frac{s}{\frac{1}{4} + \tau^2}\right),
\tag{6}
\]
再吸收 $b_0+\sum_{\tau>0}1/(\tfrac14+\tau^2)=0$，得
\[
\boxed{
\xi(s) = \frac{1}{2} \prod_{\tau > 0}
\frac{(s-\frac{1}{2})^2 + \tau^2}{\frac{1}{4} + \tau^2}
}
\tag{7}
\]
这是 (4) \textbf{在 RH 下}的简化，不是独立公理。
反之，在 $\xi$ 已由~\eqref{eq:xi-definition} 与 (4) 确定的前提下，若全部零点因子均可写成 $(s-\tfrac12)^2+\tau^2$ 之形（计重数与正规化），则 RH 成立。

\paragraph{与黎曼原稿式~(R)}
黎曼在 $s=\frac12+it$ 下写出
\begin{equation}
  \xi(t)=\xi(0)\prod_{\alpha}\Biggl(1-\frac{t^2}{\alpha^2}\Biggr),
  \label{eq:riemann-product-xi-t}
  \tag{R}
\end{equation}
其中 $\alpha$ 为 $t$-平面根的代表；式中 $\xi(0)$ 指 $s=\frac12$ 处的值，\textbf{不是} $\xi(0)=\xi(1)=\frac12$。

令 $\Xi(t):=\xi(\frac12+it)$。由偶性 $\Xi(t)=\Xi(-t)$，对 $\pm\alpha$ 配对：
\begin{equation}
  \Bigl(1-\frac{t}{\alpha}\Bigr)e^{t/\alpha}
  \cdot
  \Bigl(1+\frac{t}{\alpha}\Bigr)e^{-t/\alpha}
  =1-\frac{t^2}{\alpha^2}.
  \label{eq:pair-cancel-ch9}
\end{equation}
$t=0$ 定出常数，得
\begin{equation}
  \boxed{
  \Xi(t)
  =\xi\!\left(\frac12\right)
  \prod_{\alpha\in\mathcal{A}}
  \Biggl(1-\frac{t^2}{\alpha^2}\Biggr)
  }.
  \label{eq:Xi-riemann-form-ch9}
\end{equation}
即~(R)。该推演\textbf{不依赖 RH}。
式 (7) 是 RH 下 $s$-变量的实二次型乘积；式~(R) 是 $t$-变量对 $\pm\alpha$ 的配对，层次不同。
黎曼写出了~(R) 并简述收敛与常数，但尚未给出整函数理论下的严格证明（第~\ref{ch:riemann_original_interpretation}~章）。
下图汇总由~(4) 到~(R) 的步骤。

% 不用浮动 figure，避免排到「对数导数对照」节之后
\begin{center}
\begin{tikzpicture}[
  font=\small,
  >=Stealth,
  box/.style={
    draw=MainBlue!70!black,
    rounded corners=2pt,
    align=center,
    inner sep=6pt,
    text width=0.86\linewidth,
    fill=LightBlue!35,
  },
  arr/.style={->, very thick, MainBlue!75!black},
  steplab/.style={
    font=\footnotesize\sffamily,
    text=MainBlue!80!black,
    align=left,
    text width=0.78\linewidth,
  },
]
  \node[box] (H) {%
    \textbf{式~(4)}（Hadamard）\\[2pt]
    $\xi(s)=\xi(0)\,e^{b_0 s}
      \prod_{\rho}\bigl(1-s/\rho\bigr)e^{s/\rho}$
  };
  \node[box, below=32mm of H, fill=LightGreen!40, draw=DarkGreen!70!black] (R) {%
    \textbf{式~(R)}（黎曼，$s=\frac12+it$）\\[2pt]
    $\xi(t)=\xi(0)\prod_{\alpha}\bigl(1-t^2/\alpha^2\bigr)$
    \\[2pt]
    {\footnotesize $\xi(0)$ 指 $s=\frac12$；$\alpha$ 为 $t$-根代表}
  };
  \draw[arr] (H) -- (R)
    node[midway, right=4mm, steplab, anchor=west] {%
      (i) 换元 $s=\frac12+it$，偶性 $\Xi(t)=\Xi(-t)$\\[4pt]
      (ii) $\pm\alpha$ 配对：~\eqref{eq:pair-cancel-ch9}\\[2pt]
      \hspace{1.1em}消 $e^{\pm t/\alpha}$ 与有效 $e^{ct}$\\[4pt]
      (iii) $t=0$ 定 $c_0=\xi\!\left(\frac12\right)$
    };
\end{tikzpicture}
\captionof{figure}{由 Hadamard 乘积到黎曼 $t$-乘积}
\label{fig:hadamard-to-riemann-flow}
\end{center}

\section*{对数导数对照与本章要点}
\label{sec:log-derivative-summary}

前两节已给出 $\xi$ 的定义、Hadamard 乘积与 $\xi'/\xi$、$\zeta'/\zeta$ 的推导。
本节只作对照，便于第~\ref{ch:riemann_explicit_formula}~章查阅
（$f'/f$ 的一般性质见第~\ref{ch:complex_analysis_intro}~章第~\ref{sec:argument_principle}~节；
$-\zeta'/\zeta=\sum\Lambda(n)n^{-s}$ 见第~\ref{ch:riemann_zeta_intro}~章第~\ref{subsec:arith_dirichlet_gen}~节）。

\paragraph{已建立的公式}
\begin{align}
  \frac{\xi'}{\xi}
  &= b_0+\sum_{\rho}\Bigl(\frac{1}{s-\rho}+\frac{1}{\rho}\Bigr)
  &&\text{\eqref{eq:xi-log-derivative}}
  \\
  \frac{\xi'}{\xi}
  &= \frac{1}{s}+\frac{1}{s-1}-\frac12\log\pi
    +\frac12\frac{\Gamma'(s/2)}{\Gamma(s/2)}+\frac{\zeta'}{\zeta}
  &&\text{\eqref{eq:xi-logderiv-from-def}}
  \\
  \frac{\zeta'}{\zeta}
  &= \frac{\xi'}{\xi}-\frac{1}{s}-\frac{1}{s-1}
    +\frac12\log\pi-\frac12\frac{\Gamma'(s/2)}{\Gamma(s/2)}
  &&\text{\eqref{eq:zeta-from-xi-logderiv}}
\end{align}
第三式与~\eqref{eq:zeta-logderiv-hadamard} 在代入 Hadamard 展开后等价。

\paragraph{极点来源}
\begin{itemize}
  \item $\xi'/\xi$：极点仅在非平凡零点（$\xi$ 为整函数）；$m$ 阶零点处留数 $m$。
  \item $\zeta'/\zeta$：极点在 $\zeta$ 的零点（非平凡与平凡）及 $s=1$（留数 $-1$）。
  \item $\Gamma'/\Gamma$：极点在 $\Gamma$ 的极点；经~\eqref{eq:xi-logderiv-from-def} 进入 $\xi$ 与 $\zeta$ 的换算。
  \item $1/\zeta$：$\zeta(\rho)=0$ 处有极点；$s=1$ 处有零点。
\end{itemize}

\paragraph{选用}
零点部分分式用~\eqref{eq:xi-log-derivative} 或~\eqref{eq:zeta-logderiv-hadamard}；
算术生成函数在 $\operatorname{Re}s>1$ 用 $-\zeta'/\zeta=\sum\Lambda n^{-s}$；
$\xi$ 与 $\zeta$ 之间用~\eqref{eq:zeta-from-xi-logderiv}。
若再证明 $\zeta(1+it)\neq 0$ 并控制 $\operatorname{Re}s=1$ 附近的增长，可由此推出素数定理；
本章只建立乘积与对数导数，完整路线见第~\ref{ch:zero_distribution}、\ref{ch:riemann_explicit_formula}~章。

\begin{center}
\fbox{\parbox{0.92\textwidth}{
\centering
\vspace{0.2em}
\textbf{第十章要点}\\[6pt]
\renewcommand{\arraystretch}{1.45}
\begin{tabular}{lp{9.2cm}}
\hline
定义 & $\xi(s)=\dfrac12 s(s-1)\pi^{-s/2}\Gamma(s/2)\zeta(s)$ \\
整函数 / 对称 & 命题~\ref{prop:xi-entire}、\ref{prop:xi-fe}；$\xi(s)=\xi(1-s)$ \\
Hadamard & $\xi(s)=\xi(0)e^{b_0 s}\prod_\rho(1-s/\rho)e^{s/\rho}$，
  $b_0=\xi'(0)/\xi(0)$ \\
对数导数 & \eqref{eq:xi-log-derivative}--\eqref{eq:zeta-logderiv-hadamard} \\
原稿 $t$-乘积 & \eqref{eq:Xi-riemann-form-ch9}（$s=\frac12+it$，不依赖 RH） \\
\hline
\end{tabular}
\vspace{0.2em}
}}
\end{center}

\section*{本章小结与后续展望}\label{sec:xi_function_summary}
\addcontentsline{toc}{section}{本章小结与后续展望}

本章完成 $\xi$ 的整函数化、一阶 Hadamard 乘积，以及 $\xi'/\xi$ 与 $\zeta'/\zeta$ 的换算，
并导出与原稿衔接的 $t$-乘积。
$\xi$ 的零点即 $\zeta$ 的非平凡零点。

下一章从几何上讨论零点分布；
第~\ref{ch:prime_number_theorem}~章据此给出素数定理证明提要（$\psi\sim x$）；
第~\ref{ch:riemann_explicit_formula}~章再将 $\zeta'/\zeta$ 的零点展开代入显式公式，
连接素数计数函数的误差结构。
